\section{Derivaciones de medida y permiso local seleccionados (GCI-32)}
\label{sec:sd4-derivacion-gci32}
Se conservan dos SEL-751 y seis TC 150/5 A de GCI-29. Se seleccionan doce Phoenix URTK/S, referencia 0311087: dos bornes por secundario de cada uno de los seis TC. Cada par queda adyacente y aislado de los otros pares, con un puente de maniobra SB 2-RTK/S, referencia 0311236, seis en total, del lado TC. La ficha de URTK/S muestra el circuito de prueba con par k/l: operación normal, medida, cortocircuito de transformador y prueba de relé. Antes de separar el lado relé se cierra y comprueba el puente del lado transformador; al volver a servicio se restituye continuidad y después se retira el corto. No se unen secundarias de diferentes fases/caminos. El conjunto se monta en riel DIN dentro del tablero de medida cerrado, con marcas de polaridad y acceso exclusivo de personal autorizado. La barra de maniobra no produce supervisión automática de continuidad: el permiso se retira durante prueba y mantenimiento \parencites[pp.~1--3 y 6--10]{lote32-urtk}{lote32-sb2}.

URTK/S admite 0,5--10 mm$^2$ rígido y 0,5--6 mm$^2$ flexible, corriente nominal 41 A y tensión 400 V; acepta el conductor secundario de 1,5 mm$^2$ ya seleccionado. Se fija par de conexión 1,2--1,5 Nm y del cursor 0,6--0,8 Nm. Doce bornes de 8,2 mm ocupan 98,4 mm de riel antes de separadores/topes. Los seis puentes y los bornes reemplazan la partida genérica de seis bloques de prueba: un bloque por TC requiere dos bornes, no uno. La reserva de 0,30 VA para bornes de GCI-29 se aplica al par completo con conexiones: a 5 A equivale a $R_{\rm par}\le0,30/25=0,012~\Omega$. La ficha no publica resistencia máxima del par armado; se conserva ese valor como límite propio de aceptación de resistencia instalada, no como valor de catálogo. La carga secundaria completa debe seguir dentro de 1,7300 VA calculados y 3 VA de clase del TC; si las conexiones exceden la reserva, se recalcula antes de liberar \parencite[pp.~1--3]{lote32-urtk}.

Se seleccionan seis Schneider A9F74102, Acti9 iC60N unipolares de 2 A curva C, tres por grupo, para las derivaciones exclusivas de tensión L1/L2/L3 del SEL. Se fija conductor de cobre 1,5 mm$^2$, máximo 2 m de recorrido en el mismo tablero cerrado, segregado de potencia; neutro de medida continuo y PE del chasis independiente. No se intercala fusible en secundario TC. La variante publica Icn 6 kA a 230 V según IEC 60898-1 y capacidad de corte industrial hasta 50 kA a 220--240 V según IEC 60947-2. Esa capacidad no se atribuye a 400 V: la derivación es fase--neutro del tablero nuevo 400/230 V y se comprueba tensión real, corriente de cortocircuito y esquema de tierra. El cálculo de falla del tablero, protección del neutro según ese esquema y coordinación siguen pendientes; seleccionar 2 A no prueba selectividad ni despeje de una falla de impedancia elevada. El dispositivo se monta al origen de la derivación y se verifica apertura de cada fase en el protocolo de medida \parencite{lote32-mcb2}.

\textbf{Interfaz seleccionada.} Para el permiso SR-B se reemplaza la salida híbrida propuesta en GCI-29 por los contactos electromecánicos estándar Form A del módulo base SEL. Se conserva el relé, fuente 24--48 V DC, entradas convencionales de 5 A/300 V fase--neutro y se seleccionan entradas digitales de 24 V DC; la tarjeta adicional 4DI/4DO se pide estándar, no híbrida. El pedido configurado debe reflejar esta selección; no se inventa una cadena MOT. No se incorpora resistencia de descarga ni relé intermedio. En el manual SEL-751 de fecha 20240329, OUT101 corresponde A03--A04 y OUT102 A05--A06; se cablean en serie entre la salida del PTCB de 1 A existente de SR-B y A1 del XPS. A2/B2 conserva retorno de CEN-26. La fuente de SEL sigue detrás de su PTCB de 2 A, no detrás del contacto de permiso. Así una retirada no apaga el relé ni impide registrar la causa. Los contactos Form A abiertos por falta de alimentación o fallo interrumpen A1; no se utiliza el contacto NC de OUT103 como permiso \parencite[pp.~2.2--2.3, 2.28--2.31 y 4.200--4.201]{lote32-selmanual2024}.

La ficha actual publica para contacto estándar interrupción a 24 V DC de 0,75 A con $L/R=40$ ms. XPSUAF13AP publica 2 W DC nominales a 24 V: corriente nominal de referencia $2/24=0,08333$ A, nueve veces menor que 0,75 A. Esto identifica carga y capacidad, pero no demuestra corriente de arranque ni constante de entrada electrónica del XPS; se comprueban ambas y la caída de tensión en el circuito instalado. Los contactos estándar carecen de la especificación de fuga de estado abierto de la salida híbrida; no se le traslada el dato de 500 $\mu$A ni se declara cero absoluto de fuga por extrapolación. Si un contacto abre o pierde control, retira SR-B; un contacto soldado debe detectarse en la prueba de retirada y el otro contacto sigue una función distinta. No se atribuye independencia de procesadores, vigilancia de soldadura del SEL ni un nivel SIL al par. Un corto de cable que puentea ambos contactos puede anular la retirada y requiere segregación, inspección y verificación de circuito; no se resuelve con la lógica \parencites[p.~34]{lote29-sel751}[p.~21]{lote26-xpsmanual}.

\textbf{Ajustes de ingeniería documentados.} El manual 20240329 define OUT101FS y OUT102FS, alarmas HALARM/SALARM, elementos 50P1T, 59P1T, 27P1T, relación CTR/PTR y conexión DELTA\_Y. Con ese mapa se fija la plantilla:
\begin{quote}\ttfamily\small
CTR := 30; PTR := 1.00; DELTA\_Y := WYE; VNOM := 400;\\
FNOM := 50; PHROT := ABC;\\
50P1P := 3.75; 50P1D := 0.00; 50P1TC := 1;\\
59P1P := 239.00; 59P1D := 0.00;\\
27P1P := 215.00; 27P1D := 0.00;\\
OUT101FS := Y; OUT101 := HALARM OR SALARM;\\
OUT102FS := Y;\\
OUT102 := 50P1T OR 59P1T OR 27P1T OR LOP OR NOT IN101;
\end{quote}
IN101 recibe señal local mantenida de circuito de medida habilitado; su apertura provoca causa de retirada, no permiso. Su interfaz física/supervisión independiente queda coordinada con el permiso ambiental de INT-31; no se conecta una salida PLC arbitraria para declarar resuelto ese circuito. OUT101 mantiene la alarma de autoprueba; OUT102 usa modo fail-safe: excitado sin causa, desexcitado con causa o pérdida de control. La protección de sobretensión sigue la fase máxima y la de subtensión la mínima. Se usan los bits temporizados 27P1T y 59P1T, correspondientes respectivamente a mínimo y máximo de tensión de fase según las figuras 4.76/4.77; no se sustituyen por bits de coincidencia de las tres fases. El mapa final del pedido debe acreditar detección de una sola fase anormal. El subtensión de 215 V retira permiso por fase ausente; con tolerancia publicada $\pm1\,\%$ del ajuste más 0,5 V, el extremo inferior es 212,35 V. LOP no sustituye la prueba de fase sin tensión ni detecta cualquier fallo de TC \parencite[pp.~4.5--4.6, 4.15--4.16, 4.121--4.124, 4.145--4.146 y 4.200--4.201]{lote32-selmanual2024}.

Esta plantilla utiliza funciones y bornes reales del manual identificado, y no se presenta como fichero cargado en un equipo actual. La ficha 2026 y el manual 2024 no demuestran por sí solos identidad de firmware, tabla de opciones o bornes del pedido. Antes de energizar se cotejan manual/configuración vigentes, estados de arranque y archivo de ajustes exportado; cualquier diferencia conserva SR-B retirado. El arranque aproximado de 5--10 s del SEL no se convierte en temporizador de autorización. Recuperación de la alarma o de IN101 sólo hace disponible el rearme manual supervisado del XPS; no cierra bancada automáticamente. Se conserva el máximo de diez pasos y los espejos de los dos LC1D32P7 \parencites[p.~34]{lote29-sel751}[pp.~2.30 y 4.200--4.201]{lote32-selmanual2024}{lote29-selhealth}.

\textbf{Latencia revisada.} La salida estándar publica hasta 8 ms desde excitación de bobina hasta contacto; el modo fail-safe no obtiene apertura de 50 $\mu$s por elegir híbrido. Para esta nueva selección se retira la suma de 184 ms de GCI-29: el dato de protección menor de 1,75 ciclos está condicionado a salida híbrida y no acredita tiempo de elemento más salida estándar. Desde pérdida efectiva de 24 V en A1 de XPS, se conservan sólo 120+19=139 ms parciales para XPS/contactor. Ante una causa medida, $t_{\rm total}=t_{\rm elemento+logica+salida\ estandar}+t_{\rm XPS}+t_{\rm contacto+extincion}$: el primer término no tiene máximo completo verificado para esta configuración. La aceptación futura mide el recorrido desde inyección de causa hasta retirada de potencia y la reacción del grupo durante esa ventana; no se autoriza bancada por un límite de catálogo aún incompleto \parencites[p.~34 y pp.~38--40]{lote29-sel751}[pp.~2.30 y 4.200]{lote32-selmanual2024}[p.~45]{lote26-xpsmanual}[p.~2]{lote26-contactor}.

\section{Perfil de recarga y margen de generación (CEN/ECI-32)}
\label{sec:sd4-recarga-cen32}
La guía Eaton 93T revisión 004 distingue dos funciones: en pantalla UPS permite encender/apagar el cargador, y en especificaciones publica máximo configurable de 60 A para el modelo 60 kW. No publica en esas secciones rango inferior, paso ni tolerancia que autorice exactamente 20 A; tampoco convierte la maniobra de pantalla en una entrada de seguridad automática. Se mantiene 20 A como límite requerido de corriente real agregada por banco, no como ajuste OEM habilitado. No se asigna a CN12 ni a REPO una inhibición selectiva sin mapa de fabricante. El bloqueo local por MN se reserva a retirar entrada AC, mantiene la rama en batería y no se cicla para regular recarga \parencite[pp.~69--70 y 98]{p7-eaton-93t-guia}.

Con emisión aumentada 25\,\% y caudal útil reducido 10\,\% de INT-31, a 20 A reales y 900 m$^3$/h indicados la concentración estacionaria de cribado es 0,390833\,\%; a 21 A, 0,410375\,\%. La corriente máxima real para 0,4\,\% resulta $60(0,004)(810)/(7,5978\cdot1,25)=20,4691$ A. Por tanto sólo quedan 0,4691 A sobre veinte amperios para errores/variaciones ya no cubiertos por aquel cribado; un setpoint nominal 20 A no acredita ese límite real. La supervisión debe sumar las corrientes positivas de carga por cadena y controlar temperatura/tensión; una cadena descargando no compensa el gas de otra cargando. El SEL-751 seleccionado mide AC y no se usa como sensor de esa suma DC. No se declara seleccionado un sensor DC sin rango, aislamiento frente a 607 V, error, tiempo, fallo y salida efectiva acreditados. La autorización de recarga sigue retirada; medir la corriente total que muestra la UPS no prueba por sí sola reparto entre cadenas \parencite[pp.~97--98]{p7-eaton-93t-guia}.

Para separar generación de carga protegida se usa un cribado propio con eficiencia de diseño 0,94 previamente empleada y FP de entrada 0,99: $S_{\rm entrada}\ge(P_{\rm protegida}/0,94+P_{\rm recarga,AC}+P_{\rm auxiliares})/0,99$. En este cribado, a 60 kW y 20 A/588 V, suponiendo eficiencia ideal del cargador y cero auxiliares, $S_{\rm entrada}\ge(60/0,94+11,76)/0,99=76,3533$ kVA. Si se reserva bancada máxima superior de 11,57625 kW a 420 V, el subtotal mínimo es 87,9296 kVA antes de ventilador de bancada, control y otras entradas del grupo. Ese valor no autoriza la combinación: usa cota optimista del cargador y eficiencia de proyecto, no curva certificada. A plena demanda parcial de 72,080015 kW de CEN-26, supera además los 60 kW de la UPS vigente y el mismo cribado subiría a 100,9105 kVA con bancada. La retirada prioritaria de bancada no corrige sobrecarga de UPS ni toda demanda de servicio. Se mantienen pendientes balance por fase/estado, frío/bombas y respuesta del grupo; no se adopta el candidato 93PS ni se asigna autonomía desde estos subtotales.


Los subtotales actuales con impulsión de ENV-28 son 73,668815 kW en régimen y 75,257615 kW durante solape, antes de frío/bombas y de un margen adicional para los nuevos receptores. Bajo los mismos supuestos de cribado y con bancada, resultan respectivamente 102,6177 y 104,3250 kVA; no son cotas mínimas OEM demostradas ni demanda final medida. La incompatibilidad no se remedia manteniendo el subtotal histórico de 72,080015 kW.
